\section{二群の平均値差}
\label{s20_ttest}

\subsection{授業内容}
\subsubsection{科目の中でのこのコマの位置づけ}
\begin{tabular}[t]{cccccc}
    記述統計[3] & 確率[6] & モデル[9] & 推測・推定[2] & 検定[4] & 実践[4]
\end{tabular}

\subsubsection{概要}
心理学の実践で最もよく用いられるのが，平均値差の検定である。帰無仮説検定としては，t検定を行うことになる。対応のあるt検定，対応のないt検定の順に，Rでの演習も交えて伝える。その際，結果から適切な読み取りができるようになること，記述することなどが求められる。統計環境はRを用いるが，SPSSそのほかの統計パッケージを用いても同様の出力が得られる。統計環境が変わる可能性もあるので，出力画面の位置で内容を覚えるのではなく，意味的な理解をしておく必要がある。
\subsubsection{コマ主題細目(箇条書き)}
\begin{description}

\item[一要因Withinデザイン，2水準]一要因Wtihinデザイン，2水準モデルの検定統計量もまた，$t$分布を用いる。これは「対応がある」，すなわち個人が特定されているため，正規分布する変数の差の分布についての検定であり，既に習った1つの平均の差についての検定と同じである。
\begin{flushright}
    $\to$ \citet{橋本201607}のP.1--13.
\end{flushright}

\item[一要因Betweenデザイン，2水準]一要因Betweenデザイン，2水準モデルの検定統計量は，$t$分布を用いるため$t$検定と呼ばれることもある。帰無仮説は二群の平均値差がないとするものであり，5\%水準による検定結果と差の信用区間が示されることが多い。また効果量$d$を別途計算して報告されることにも注意する。
\begin{flushright}
    $\to$　\citet{橋本201607}のP.14--25.
\end{flushright}

% \item[乱数を用いたシミュレーション]乱数を用いることで，具体的に帰無仮説の想定する空間やサンプリングの意味，検定結果の現れ方が分かりやすくなる。t検定を例に，2群の平均値の差，母SDの大きさなどを変え，検定がどのような傾向を持っているかを知る。同時に，平均値差を標準偏差で割った効果量のことを考える。
 
% \item[基本的操作のおさらい]まずは基本操作を復讐しておく。データフレーム，変数の呼び出しに加え，確率分布関数の呼び出し方，乱数の使い方などを再確認する。

\item[Rによる実践；一要因Withinデザイン，2水準]このデザインはRではt.test関数を用いる。Withinデザインの場合はpairedオプションをつける必要がある。効果量の算出には，別途\texttt{effsize}パッケージを導入すると良い。
\begin{flushright}
    $\to$ \citet{橋本201607}のP.1--25.
\end{flushright}

\item[Rによる実践；一要因Betweenデザイン，2水準]このデザインはRではt.test関数を用いる。Betweenデザインの場合はpairedオプションをつける必要はないが，Welchの補正がかけられている。仮説検定の前提となった仮定を思い出し，これが満たされない場合はこうした補正がかかることを学ぶ。
\begin{flushright}
    $\to$ \citet{橋本201607}のP.1--25.
\end{flushright}

\end{description}
\subsubsection{キーワード}
\begin{itemize}
    \item 対応のないt検定，対応のあるt検定
    \item 線形モデルと検定
    \item t検定
    \item 効果量
\end{itemize}
\subsection{授業情報}
\paragraph{コマの展開方法}
講義および演習
\paragraph{標準シラバスにおける位置づけ}
科目番号5；心理学統計法；1.心理学で用いられる統計手法；H 推測統計:代表値と散布度をめぐって(2)
科目番号5 心理学統計法；2統計に関する基礎的な知識；C/D/E エクセル，R，SPSS入門
\subsubsection{予習・復習課題}
\paragraph{予習}
RやRStudio，プロジェクトによる管理など基本的なことを改めて教示しないので，これまでの復習をかねて基本的なR/RStudioの挙動について再確認しておくこと。データファイルは事前にCoursePowerなどを通じて配布されるが，読み込みなどについては逐一立ち止まって解説しないので，できることを確認して授業に挑むこと。
\paragraph{復習}
時間内に収まりきらなかった課題については，復習用の課題となり，Rmdファイルで提出することが求められる。適切な分析ができていない場合は，課題を不受理として扱う。これまでと同様，課題が提出されていることが単位取得の要件になることに注意し，しっかり対応して欲しい。
また，今回の内容について\citet{橋本201607}が丁寧な解説をしているため，十分理解できなかった場合は同書を一読して理解しておく。内容的には，\citet{Yamada2004}のP.162--207も該当する。最も，これらを実践的にRで行うことが目的であり，Rの使い方に言及している\citet{kosugi2019}なども参考にしてほしい。
効果量に関しては，\citet{大久保201201}のP.52--68を参照しておく。